% p.26 — الباب الرابع العشر في دعا يدعو (ناقص واوي)
\mspage{26}
\bab{\com{الباب الرابع العشر}{كذا في الأصل؛ ولعل الصواب: الباب الرابع عشر} في دعا يدعو (ناقص واوي)}
% Layout: six tall columns (right-to-left), each a stack of boxed blocks of unequal height.
% Column 1 is the summary column: the باب name (نَصَرَ) in its own box, then the ten
% underlying (pre-iʿlāl) forms دَعَوَ … لَا تَدْعُوْ, then the ten actual forms دَعَا … لَا تَدْعُ.
% In columns 2–6 every sub-heading box holds the heading and, beneath it, the underlying form
% (e.g. «ماضي معروف / دَعَوَ»); the regular paradigm follows outside the box.
% Row alignment is approximate (all columns run to 21 lines here); \cline marks the printed box
% boundaries inside each column. Forms ending in alif maqṣūra carry a small (dagger) alif
% in the lithograph (يُدْعٰى); it is transcribed as U+0670.
\begin{sarf}{|c|c|c|c|c|c|}
نَصَرَ & \thead{ماضي معروف} & تَدْعُوْ & دُعِيَ & يُدْعَيْنَ & اُدْعُوْا\\\cline{1-1}
دَعَوَ & دَعَوَ & تَدْعُوَانِ & دُعِيَا & تُدْعٰى & اُدْعِيْ\\\cline{2-2}
يَدْعُوُ & دَعَا & يَدْعُوْنَ & دُعُوْا & تُدْعَيَانِ & اُدْعُوَا\\
دَعْوًا & دَعَوَا & تَدْعُوْ & دُعِيَتْ & تُدْعَوْنَ & اُدْعُوْنَ\\\cline{6-6}
دَاعِوٌ & دَعَوْا & تَدْعُوَانِ & دُعِيَتَا & تُدْعَيْنَ & \thead{نهي حاضر}\\
دُعِوَ & دَعَتْ & تَدْعُوْنَ & دُعِيْنَ & تُدْعَيَانِ & لَا تَدْعُوْ\\\cline{6-6}
يُدْعَوُ & دَعَتَا & تَدْعِيْنَ & دُعِيْتَ & تُدْعَيْنَ & لَا تَدْعُ\\
دَعْوًا & دَعَوْنَ & تَدْعُوَانِ & دُعِيْتُمَا & أُدْعٰى & لَا تَدْعُوَا\\
مَدْعُوْوٌ & دَعَوْتَ & تَدْعُوْنَ & دُعِيْتُمْ & نُدْعٰى & لَا تَدْعُوْا\\\cline{5-5}
اُدْعُوْ & دَعَوْتُمَا & أَدْعُوْ & دُعِيْتِ & \thead{اسم مفعول} & لَا تَدْعِيْ\\
لَا تَدْعُوْ & دَعَوْتُمْ & نَدْعُوْ & دُعِيْتُمَا & مَدْعُوْوٌ & لَا تَدْعُوَا\\\cline{1-1}\cline{3-3}\cline{5-5}
دَعَا & دَعَوْتِ & \thead{اسم فاعل} & دُعِيْتُنَّ & مَدْعُوٌّ & لَا تَدْعُوْنَ\\\cline{6-6}
يَدْعُوْ & دَعَوْتُمَا & دَاعِوٌ & دُعِيْتُ & مَدْعُوَّانِ & \thead{مضارع بلم}\\\cline{3-3}
دَعْوَةً & دَعَوْتُنَّ & دَاعٍ & دُعِيْنَا & مَدْعُوُّوْنَ & لَمْ يَدْعُوْ\\\cline{4-4}\cline{6-6}
دَاعٍ & دَعَوْتُ & دَاعِيَانِ & \thead{مضارع مجهول} & مَدْعُوَّةٌ & لَمْ يَدْعُ\\
دُعِيَ & دَعَوْنَا & دَاعُوْنَ & يُدْعَوُ & مَدْعُوَّتَانِ & لَمْ يَدْعُوَا\\\cline{2-2}\cline{4-4}
يُدْعٰى & \thead{مضارع معروف} & دَاعِيَةٌ & يُدْعٰى & مَدْعُوَّاتٌ & لَمْ يَدْعُوْا\\\cline{5-5}
دَعْوَةً & يَدْعُوُ & دَاعِيَتَانِ & يُدْعَيَانِ & \thead{أمر حاضر} & لَمْ تَدْعُ\\\cline{2-2}
مَدْعُوٌّ & يَدْعُوْ & دَاعِيَاتٌ & يُدْعَوْنَ & اُدْعُوْ & لَمْ تَدْعُوَا\\\cline{3-3}\cline{5-5}
اُدْعُ & يَدْعُوَانِ & \thead{ماضي مجهول} & تُدْعٰى & اُدْعُ & لَمْ يَدْعُوْنَ\\
لَا تَدْعُ & يَدْعُوْنَ & دُعِوَ & تُدْعَيَانِ & اُدْعُوَا & لَمْ تَدْعُ\\
\end{sarf}
